\subsection{Validation Design}
\label{subsec:validation}

The framework claims to handle environmental disruption regardless of how the evidence is represented. That claim cannot be tested with a single city, because a single city exercises one data type and leaves the rest unverified. We therefore designed the two case studies as complementary validation instances, each chosen to exercise a distinct capability of the mechanism described in Section~\ref{sec:penalization}.

\subsubsection{Why Two Cases and Not One}

A validation instance is useful only if it can fail. A city with a single environmental layer would confirm that the pipeline runs, not that it generalizes, since the same result could be produced by code specialized to that layer type. The two cases are therefore selected to differ along the dimensions where the mechanism makes commitments.

\textbf{Continuity of the effect.} Curitiba is modeled with land surface temperature, a continuous raster field. The disruption there raises travel cost in proportion to a measured quantity, and no connection is ever removed: every path exists in both scenarios, with different costs. \textbf{Discreteness of the effect.} Porto Alegre is modeled with geological risk polygons, a categorical vector layer. The disruption there renders segments impassable, and some origins lose access entirely.

Between them, the two cases cover both branches of the modifier mechanism — continuous sampling and geometric intersection — and both outcomes of the cost function: penalization and interdiction. Neither city alone does so. A framework validated only in Curitiba could not claim to support interdiction; validated only in Porto Alegre, it could not claim to support continuous penalization.

\subsubsection{What Each Case Validates}

\textbf{Curitiba validates sampling and aggregation.} Because the thermal layer is a raster, the cost of an edge depends on which pixels the edge crosses. The case exercises the choice between midpoint and segment sampling, and the difference between them is observable: edges crossing a thermal boundary yield different conditioned costs under the two strategies. Curitiba also exercises the composite arrangement, in which vegetation attenuates the thermal penalty only where the same location is hot, which requires that both layers be read at identical sample points. A defect in that correspondence would appear as attenuation in wrong places, and the case would expose it.

\textbf{Porto Alegre validates partition and interdiction.} Because the risk layer is a polygon set, the cost of an edge depends on how its geometry is split at polygon boundaries. The case exercises the partitioning procedure, the treatment of overlapping polygons, and the propagation of infinite cost through routing. Porto Alegre also validates the distinction between attenuation and interdiction at the level of the indicator: the framework must report cells that lost all access separately from cells whose access merely became slower, since only the first is an absence of access.

\subsubsection{What the Pair Validates Together}

The two cases share the same core, the same configuration schema, and the same indicator computation. Applying them without code change is itself the evidence for the generality claimed in the Introduction: the difference between the experiments lies entirely in the data and the declared rules, not in the implementation.

This is the sense in which the case studies are a validation instrument rather than an illustration. Each one is constructed to stress a specific commitment of the mechanism, and the pair is constructed so that no commitment is left untested. Section~\ref{sec:case_studies} describes the data and the calibration of each city; Section~\ref{sec:results} reports what the pair produced.
